(1) 边界向恒外压移动，对外功 \(W_{\rm irr}=p_{\rm ext}(V_f-V_i)=nRT(1-V_i/V_f)\)，为正。理想气体等温初末态内能相等，\(\Delta U=0\)，第一定律 \(\Delta U=Q-W\) 给 \(Q=W_{\rm irr}\)，由热浴供应。
(2) 气体熵是状态函数，\(\Delta S_{\rm gas}=nR\ln(V_f/V_i)\)；恒温热浴失热 \(Q\)，故 \(\Delta S_{\rm bath}=-Q/T=-nR(1-V_i/V_f)\)。因此总熵产生 \(\Delta S_{\rm gen}=nR[\ln z-1+1/z]\)，其中 \(z=V_f/V_i>1\)。
(3) 定义 \(f(z)=\ln z-1+1/z\)，则 \(f(1)=0\)，\(f'(z)=1/z-1/z^2=(z-1)/z^2>0\) 对 \(z>1\) 成立，故 \(\Delta S_{\rm gen}>0\)。同端点的可逆等温功为 \(W_{\rm rev}=nRT\ln z\)，因为 \(\ln z>1-1/z\)，可逆功严格更大。
(4) 当 \(z\to1\)，展开 \(f(z)=\frac12(z-1)^2+O((z-1)^3)\)，熵产生趋零；当 \(z\gg1\)，实际外功趋 \(nRT\)，而可逆功随 \(\ln z\) 增大。不可逆中间态不保持均匀的准静态热力学压强，所以真实边界功必须使用外压 \(p_{\rm ext}\)，不能直接沿可逆状态路径积分系统压强。
